\documentclass{article}
\usepackage[UTF8]{ctex}
\usepackage{graphicx}
\usepackage[top=3cm, bottom=2cm, left=2.5cm, right=2.5cm]{geometry}
\usepackage{array}
\usepackage{enumitem}
\usepackage{multirow}
\usepackage{xcolor}
\usepackage{float}
\usepackage{titlesec}
\usepackage{fancyhdr}
\usepackage{hyperref}
\definecolor{oucblue}{RGB}{0,84,158}
\hypersetup{
    colorlinks=true,
    linkcolor=blue,
    citecolor=blue,
    urlcolor=blue,
    filecolor=blue
}
\pagestyle{fancy}
\fancyhf{}
\fancyhead[L]{\includegraphics[height=0.8cm]{ouc.png}}
\fancyhead[R]{\nouppercase{\leftmark}}
\fancyfoot[C]{\thepage}
\renewcommand{\headrulewidth}{0.4pt}
\renewcommand{\footrulewidth}{0pt}

\title{实验报告3}
\author{伍紫涵}


\titleformat{\section}[block]
{\normalfont\Large\bfseries\centering}{\thesection}{1em}{}

\begin{document}


\begin{titlepage}
    \centering
    {\Huge\bfseries 实验报告3 \par}
    \vspace{2cm}
    \includegraphics[width=0.7\linewidth]{ouc.png} \par
    \vspace{0.6cm}
    {\zihao{2}\bfseries\color{oucblue} 中国海洋大学}\par
    \vspace{0.5cm}
    {\LARGE 姓名：伍紫涵 \par}
    \vspace{0.4cm}
    {\LARGE 专业：计算机科学与技术 \par}
    \vspace{0.4cm}
    {\large \today \par}
\end{titlepage}

\thispagestyle{plain}
\begin{abstract}
{\large
本文是我学习的关于命令行环境、python入门基础和python视觉应用的实验报告，接下来我会详细地介绍我学习的内容，本文的LaTeX代码已经上传到了\href{https://github.com/qiqiqisi/Latex_Used.git}{github仓库}中，可以点击超链接进行查看。
}
\end{abstract}

\pagenumbering{roman}
\setcounter{page}{1}
{
\hypersetup{hidelinks}
\tableofcontents
}


\newpage
\pagenumbering{arabic}
\setcounter{page}{1}

%---------------------------------------------------------%

\section{命令行环境}
\subsection{进程的强制停止}
进程正在进行的时候，按下Ctrl + C的时候会释放SIGINT信号，但是当代码捕获了SIGINT之后就无法停止，使用Ctrl + \\可以强制停止，它释放的是SIGQUIT。
\begin{figure}[H]
    \centering
    \includegraphics[width=0.8\linewidth]{code_show.png}
\end{figure}
\begin{figure}[H]
    \centering
    \includegraphics[width=0.8\linewidth]{code_run.png}
\end{figure}


\subsection{进程的暂时停止与继续运行}
进程正在进行的时候，按下Ctrl + Z的时候就会中止运行，输入fg的时候进程不会从头来过，而是会在上一次的基础上继续运行。
\begin{figure}[H]
    \centering
    \includegraphics[width=1\linewidth]{continue.png}
\end{figure}


\subsection{让进程在后台执行}
进程中止后，输入bg的时候进程会进入后台继续运行。
\begin{figure}[H]
    \centering
    \includegraphics[width=1\linewidth]{bg.png}
\end{figure}

\newpage
\subsection{开辟工作台运行程序}
输入tmux new -s workspace可以创建工作台，在这个工作台中可以运行程序。
\begin{figure}[H]
    \centering
    \includegraphics[width=1\linewidth]{tmux0.png}
\end{figure}
\begin{figure}[H]
    \centering
    \includegraphics[width=1\linewidth]{tmux.png}
\end{figure}

\newpage
\subsection{在工作台中分割页面}
进入工作台后，按下Ctrl + B，然后如果打出“就可以创建在竖直方向上切割窗口，如果打出的是\%，就会在水平方向上切割窗口。
\begin{figure}[H]
    \centering
    \includegraphics[width=1\linewidth]{split.png}
\end{figure}


\subsection{创建和删除别名}
使用alias可以创建别名，使用unalias可以删除创建了的别名。
\begin{figure}[H]
    \centering
    \includegraphics[width=1\linewidth]{alias.png}
\end{figure}


\section{python入门基础}
\subsection{打印语句}
使用print可以打印语句。
\begin{figure}  [H]
    \centering
    \includegraphics[width=1\linewidth]{print.png}
\end{figure}


\subsection{for循环语句}
使用for执行循环语句。
\begin{figure}[H]
    \centering
    \includegraphics[width=1\linewidth]{for.png}
\end{figure}


\newpage
\subsection{while循环语句}
使用while执行循环语句。
\begin{figure}[H]
    \centering
    \includegraphics[width=1\linewidth]{while.png}
\end{figure}


\newpage
\subsection{if分支语句}
使用if执行分支语句。
\begin{figure}[H]
    \centering
    \includegraphics[width=1\linewidth]{if.png}
\end{figure}


\newpage
\subsection{在列表中使用for语句}
在列表中使用for语句，减少了行数。
\begin{figure}[H]
    \centering
    \includegraphics[width=1\linewidth]{for_list.png}
\end{figure}


\newpage
\subsection{在列表的末尾插入数据}
使用append在列表的末尾插入数据。
\begin{figure}[H]
    \centering
    \includegraphics[width=1\linewidth]{apend.png}
\end{figure}
\newpage


\subsection{列表自动排序}
使用sort自动让列表的数据排序。
\begin{figure}[H]
    \centering
    \includegraphics[width=1\linewidth]{sort.png}
\end{figure}


\newpage
\subsection{判断容器中是否存在某个元素}
in可以直接判断某个元素是否在容器中。
\begin{figure}[H]
    \centering
    \includegraphics[width=1\linewidth]{in.png}
\end{figure}


\newpage
\subsection{反转列表中的元素}
reverse可以反转列表中的元素。
\begin{figure}[H]
    \centering
    \includegraphics[width=1\linewidth]{reverse.png}
\end{figure}


\newpage
\subsection{获取容器长度}
len可以获取容器的长度，即是容器中的元素个数。
\begin{figure}[H]
    \centering
    \includegraphics[width=1\linewidth]{len.png}
\end{figure}


\newpage
\section{python视觉入门}
\subsection{调整图片的大小}
使用resize重新调整图片的长和宽。
\begin{figure}[H]
    \centering
    \includegraphics[width=1\linewidth]{resize_code.png}
\end{figure}
\begin{figure}[H]
    \centering
    \includegraphics[width=1\linewidth]{resize_pic.png}
\end{figure}


\newpage
\subsection{增强对比度}
使用Contrast增强对比度。
\begin{figure}[H]
    \centering
    \includegraphics[width=1\linewidth]{contrast_code.png}
\end{figure}
\begin{figure}[H]
    \centering
    \includegraphics[width=1\linewidth]{contrast_pic.png}
\end{figure}


\newpage
\subsection{增强亮度}
使用Brightness增强亮度。
\begin{figure}[H]
    \centering
    \includegraphics[width=1\linewidth]{brightness_code.png}
\end{figure}
\begin{figure}[H]
    \centering
    \includegraphics[width=1\linewidth]{Brightness_pic.png}
\end{figure}


\subsection{增强锐度}
\begin{figure}[H]
    \centering
    \includegraphics[width=1\linewidth]{sharpness.png}
\end{figure}
\begin{figure}[H]
    \centering
    \includegraphics[width=1\linewidth]{sharpness_pic.png}
\end{figure}
\end{document}



